\begin{abstract}
Air traffic management inefficiencies within terminal maneuvering areas significantly increase environmental and operational costs for the aviation industry. Predicting the additional fuel burn resulting from these inefficiencies remains a complex challenge. In this study, we propose a machine learning framework designed to estimate the excess fuel consumption, recognized internationally as Key Performance Indicator 16, for commercial aircraft arriving at Guarulhos International Airport in Brazil. Our methodology integrates real-world radar trajectories, aircraft performance parameters, and additional transit time metrics into a robust data architecture. We evaluated multiple gradient boosting algorithms to capture the complex relationships between flight trajectory variables and fuel consumption. The comparative analysis reveals that, after hyperparameter optimization, the CatBoost model achieves the most effective balance of predictive performance and generalization, yielding a coefficient of determination of 0.6259 and a root mean squared error of 33.56 kilograms. Furthermore, our feature importance analysis indicates that the predicted transit time within the terminal area is the most critical explanatory variable, emphasizing the direct impact of air traffic congestion and sequencing inefficiencies on fuel burn. Ultimately, the developed approach provides stakeholders with a reliable, data-driven tool for the continuous monitoring of operational efficiency. This framework supports air navigation service providers in optimizing approach procedures, enhancing governance, and aligning operations with global aviation sustainability goals.
\end{abstract}